\documentclass[10pt,a4paper]{amsart}
\usepackage[margin=19mm]{geometry}
\usepackage{amsmath,amssymb,mathtools}
\usepackage[scr=rsfs,cal=euler]{mathalfa}
\usepackage{xfrac}
\usepackage[unicode,hidelinks,bookmarks=false]{hyperref}
\newcommand{\ie}{i.e.\ }
\newcommand{\cf}{cf.\ }
\newcommand{\resp}{respectively\ }
\newcommand{\Subsection}{\subsection}
\providecommand{\nobreakdash}{\nobreak\hbox{-}}
\setlength{\emergencystretch}{2em}
\multlinegap0pt
\usepackage{bm}
\usepackage{bbm}
\usepackage{dsfont}
\usepackage{xspace}
\usepackage{cancel}
\usepackage{mathtools}
\usepackage{cleveref}
\usepackage{amsthm}
\makeatletter
\@ifundefined{assumption}{\newtheorem{assumption}{Assumption}}{}
\@ifundefined{theorem}{\newtheorem{theorem}{Theorem}}{}
\makeatother
\makeatletter
\def\EE{\mathbb{E}}
\def\KK{\mathbb{K}}
\def\RR{\mathbb{R}}
\def\XX{\mathbb{X}}
\def\YY{\mathbb{Y}}
\def\cE{\mathcal E}
\def\cH{\mathcal H}
\def\cR{\mathcal R}
\def\cX{\mathcal{X}}
\def\cZ{\mathcal Z}
\expandafter\ifx\csname carlist\endcsname\relax
\newenvironment{carlist}
{\begin{list}{$\bullet$}
		{\setlength{\topsep}{0.1in} \setlength{\partopsep}{0in}
			\setlength{\parsep}{0.1in} \setlength{\itemsep}{\parskip}
			\setlength{\leftmargin}{0.15in} \setlength{\rightmargin}{0.08in}
			\setlength{\listparindent}{0in} \setlength{\labelwidth}{0.08in}
			\setlength{\labelsep}{0.1in} \setlength{\itemindent}{0in}}}
\fi
\def\l{\left}
\def\r{\right}
\def\tr{{\rm Tr}}
\def\two{\text{two}}
\def\wh{\widehat}
\def\wt{\widetilde}
\makeatother
\title[Adaptive Kernel Ridge Regression with Linear Structure: Shar]{Adaptive Kernel Ridge Regression with Linear Structure: Sharp Oracle Inequalities and Minimax Optimality}
\author{Xin Bing, Chao Wang}
\date{Source: arXiv:2605.11806v1 (CC BY 4.0)}
\begin{document}
\maketitle

\noindent\emph{Source and licence.} Adaptive Kernel Ridge Regression with Linear Structure: Sharp Oracle Inequalities and Minimax Optimality. Version arXiv:2605.11806v1 (CC BY 4.0), distributed under the Creative Commons Attribution 4.0 International licence, as recorded in the arXiv licence metadata for this exact version and shown on the abstract page https://arxiv.org/abs/2605.11806v1. Full rights notice: licenses/arxiv-2605.11806-CC-BY-4.0.txt.\par\smallskip

\noindent\textit{Editorial scope.} Counted unit: the Theorem whose source label is \texttt{thm\_twostep} in the article (source0.tex lines 2436--2446, source label \texttt{thm\_twostep}), delivered together with the article's own complete original proof of it (source0.tex lines 2482--2531, one \texttt{proof} environment of 50 lines). No mathematics is written, repaired or re-derived: statement and proof are the article's own text line for line. The proof is detached from the statement in the source: the article states the result at lines 2436--2446 and prints its proof at lines 2482--2531, and the proof environment's own header names the statement, so the association is the article's own. Required context retained uncounted: model (lines 1317--1319); ass\_pd\_K (lines 1485--1487); thm\_krr\_fix (lines 1496--1502); thm\_fix (lines 1525--1531). Every definition, display and label of those lines is retained unchanged. Omitted from this package and not redistributed: 1--1316, 1320--1484, 1488--1495, 1503--1524, 1532--2435, 2447--2481, 2532--2951. Nothing inside a retained excerpt is omitted or altered: every retained body is delivered exactly as the article's lines are written, so no omission receipt of any kind accompanies this package. Citations: the retained excerpts carry no citation command at all, so no bibliography entry of the article is transcribed and no reference is redistributed. Reference adaptation: none was needed and none was made; every reference inside the delivered unit resolves inside it, so no unresolved reference or citation is produced. Printed slips kept literal: no printed word, symbol, hypothesis or number of the counted statement or of its proof was altered. Layout and class: the article's own class is \texttt{article} with packages including amsmath, bm, amsthm, hyperref, tikz, bookmark, amssymb, amsfonts, titlesec, verbatim, fullpage, color, xcolor, mathrsfs; the wrapper is the lane's pinned \texttt{amsart} editorial wrapper at 10pt on A4, which restates the theorem-like environments the retained text needs and adds the article's own macro definitions that the retained text needs (\texttt{\textbackslash EE}, \texttt{\textbackslash KK}, \texttt{\textbackslash RR}, \texttt{\textbackslash XX}, \texttt{\textbackslash YY}, \texttt{\textbackslash cE}, \texttt{\textbackslash cH}, \texttt{\textbackslash cR}, \texttt{\textbackslash cX}, \texttt{\textbackslash cZ}, \texttt{\textbackslash l}, \texttt{\textbackslash r}, \texttt{\textbackslash tr}, \texttt{\textbackslash two}, \texttt{\textbackslash wh}, \texttt{\textbackslash wt}), with extra wrapper packages (bm, bbm, dsfont, xspace, cancel, mathtools, cleveref). The delivered numbering is the unit's own. Assets: no retained span contains a figure environment, a graphics-inclusion command, a PDF-page inclusion, a table or a TikZ picture; no image file is copied into this package, the package declares no asset dependency and no image file is redistributed with it. Licence: version arXiv:2605.11806v1 of this article is distributed under the Creative Commons Attribution 4.0 International licence, as recorded in the arXiv licence metadata for this exact version and shown on the abstract page https://arxiv.org/abs/2605.11806v1. A full rights notice is delivered at licenses/arxiv-2605.11806-CC-BY-4.0.txt. Characters: every retained excerpt is pure ASCII, so the delivered engine, XeLaTeX with its default Latin Modern text font, reports no missing character.
\begin{equation}\label{model}
   Y_i=f^*(X_i)+\epsilon_i,
\end{equation}

\begin{assumption}\label{ass_pd_K}
    The kernel function $K:{ \cX \times \cX} \to \RR$ is   symmetric and positive semi-definite.\footnote{We say $K$ is positive semi-definite if for all finite sets $\{x_1,\ldots,x_m\}\subset \cZ$ the $m\times m$ matrix  whose $(i,j)$ entry is $K(x_i, x_j)$ is positive semi-definite.} 
\end{assumption}

\begin{theorem}\label{thm_krr_fix}
    Grant model (\ref{model}) and
    Assumption   \ref{ass_pd_K}.  For any $\lambda >0$, one has 
    \begin{equation}\label{rate_krr_fix}
        \cR(\wt g-f^*) ~ \le ~    \inf_{f\in \cH_K}\left\{  \cR(f-f^*) + \lambda \|f\|_K^2  \right\}+ {\sigma^2  \over n}  \sum_{i=1}^n\left(\wh \mu_i\over \wh \mu_i + \lambda\right)^2.
    \end{equation}
\end{theorem}

\begin{theorem}\label{thm_fix}
Grant model (\ref{model}) and
Assumption  \ref{ass_pd_K}. For any $\lambda >0$, one has
\begin{align}\label{bd_pred_rate_fix}
      \cR(\wh f-f^*) \le   \inf_{\alpha \in \RR^d, f\in \cH_K}\left\{  \cR\l(\langle \alpha,\rangle + f-f^*\r) +  \lambda \|f\|_K^2  \right\}+ {\sigma^2\over n}  \left(  d + \sum_{i=1}^n \left(\wh\nu_i\over \wh\nu_i + \lambda\right)^2\right).
\end{align}    
\end{theorem}

\begin{theorem}\label{thm_twostep}
    Grant model (\ref{model}) and
    Assumption   \ref{ass_pd_K}.  For any $\lambda >0$, one has 
    \begin{equation}
        \label{risk_bd_two}
        \begin{aligned}
      \cR(\wh f_{\rm two}-f^*) ~ \le ~    & \inf_{\alpha\in \RR^d, f\in \cH_K}\left\{  \cR(\langle \alpha,\rangle + f- f^*)  +  {1\over n} \|P_{\XX} f(\XX)\|^2+  \lambda \|f\|_K^2  \right\}\\
      & ~~ + {\sigma^2  \over n}\left\{ d +  \tr\l(Q_{\XX} \KK^2 \KK_{\lambda}^{-2} \r)\right\}
\end{aligned}
    \end{equation}
\end{theorem} 

\begin{proof}[Proof of \cref{thm_twostep}]
  We follow similar arguments in the proof of \cref{thm_fix}.   Start with
  \begin{align*}
 \wh f_{\two} (\XX) - f^*  (\XX ) 
&
~ =~ P_ \XX  \YY  +  \KK\KK_{\lambda} ^{-1} Q_{\XX} \YY -  f^*  (\XX) 
\\
&
~ =~ P_ \XX  \cE   +   \KK\KK_{\lambda} ^{-1} Q_{\XX} \cE +  \KK\KK_{\lambda}^{-1}   Q_{\XX}  f^*  (\XX)    - Q_{\XX} f^*  (\XX) \\
&
~ =~ P_ \XX  \cE   +   \KK\KK_{\lambda}^{-1}  Q_{\XX} \cE -  \lambda  \KK_{\lambda}^{-1} Q_{\XX}  f^*  (\XX) 
\end{align*}
Therefore, the prediction risk of the predictor $\wh f_{\two}$  can be decomposed into   
\[
\cR(\wh f_{\two}) ~ =~ 
 {1 \over n} \EE \l\| P_ \XX  \cE   +   \KK\KK_{\lambda} ^{-1}  Q_{\XX} \cE\r\| ^2 +  {1 \over n}\l \|  \lambda \KK_{\lambda}^{-1}  Q_{\XX}  f^*  (\XX)\r\|^2 .
\]
Further note that
\begin{equation}
    \label{decom_two}
    \begin{aligned}
 {1 \over n} \EE \| P_ \XX   \cE   + \KK\KK_{\lambda}  Q_{\XX} \cE\| ^2  
  & =   {\sigma^2  \over n} \left\{ d +    \tr\l(\KK\KK_{\lambda} ^{-1}Q_{\XX}  \KK\KK_{\lambda} ^{-1}\r) +2 \tr\l(P_{\XX}\KK\KK_{\lambda} ^{-1}Q_{\XX}  \r)\right\}
   \\
   & ={\sigma^2  \over n} \left\{ d +     \tr\l( Q_{\XX} \KK^2 \KK_{\lambda}^{-2} \r)  \right\}
   ,
\end{aligned}
\end{equation}
where  the last step follows from 
\begin{align*}
    \tr\l(P_{\XX} \KK \KK_{\lambda} ^{-1}Q_{\XX}  \r)     & =    \tr\l(P_{\XX}\KK\KK_{\lambda} ^{-1} \r)- \tr\l(P_{\XX}\KK\KK_{\lambda} ^{-1}  P_{\XX} \r) \\
    & = \tr\l(P_{\XX}\KK\KK_{\lambda} ^{-1} \r)- \tr\l(P_{\XX}\KK\KK_{\lambda} ^{-1} \r) = 0. 
\end{align*}

To bound the second term in the right-hand of \eqref{decom_two}, we follow the same argument of proving \cref{thm_krr_fix} by replacing $f^*(\XX)$ by $Q_{\XX} f^*(\XX)$ to obtain
\begin{align*}
     {1 \over n}\l \|  \lambda \KK_{\lambda}^{-1}  Q_{\XX}  f^*  (\XX)\r\|^2 ~ \le~   \inf_{f\in \cH_K}\l\{
            {1\over n} \|f(\XX) -Q_{\XX} f^*(\XX)\|^2 + \lambda \|f\|_K^2
        \r\}.
\end{align*}
By noting that
\begin{align*}
   \|f(\XX) -Q_{\XX} f^*(\XX)\|^2   &  
   = \|P_{\XX} f(\XX) +  Q_{\XX} f(\XX) -Q_{\XX} f^*(\XX)\|^2 \\
   & = \|P_{\XX} f(\XX) \|^2+ \|  Q_{\XX} f(\XX) -Q_{\XX} f^*(\XX)\|^2
   \\
   & = \|P_{\XX} f(\XX) \|^2+ \inf_{\alpha\in \RR^d} ~\l \| \XX\alpha+  f(\XX) - f^*(\XX)\r\|^2, 
\end{align*}
we complete the proof. 
\end{proof}

\end{document}
